\kafkasection{Derive the pressure, density, and temperature profiles of an adiabatically stratified plane-parallel atmosphere under constant gravitational acceleration $g$. 
Assume that the atmosphere consists of an ideal gas of molecular weight $\mu$.
Assume that the Earth’s atmosphere is adiabatically stratified and plane parallel and that it consists of an ideal gas of molecular weight $\mu = 14$.
If the gravitational acceleration is $g = 981 \text{cm s}^{-2}$ and the temperature at sea level is $300 K$, what is the temperature at the summit of Mt. Everest at a height of $8848$ m? 
What is the pressure on the summit compared to the pressure at sea level?}
Let's start with our gas in hydrostatic equilibrium.
\begin{equation}
    \nabla P = -\rho\nabla\phi
\end{equation}
Because we have planar symmetry of our gravitational potential, we get:
\begin{equation}
    \odv{P}{z} = -\rho g
\end{equation}
The expression for $\rho$ can be replaced by $\rho$ from the polytropic equation of state:
\begin{equation}
    P = K \rho^\Gamma
\end{equation}
\begin{equation}
    \odv{P}{z} = -g\ab(\frac{P}{K})^{1/\Gamma}
\end{equation}
This is a separable equation so let's solve for $P$:
\begin{equation}
    P^{-1/\Gamma} \d P = -K^{-1/\Gamma} g\d z
\end{equation}
For the position z, we integrate from $0\rightarrow z$ and the pressure we integrate from $P_0\rightarrow P(z)$ (pressure at the surface to pressure at a given z):
\begin{equation}
    \ab(\frac{\Gamma}{\Gamma - 1})\ab(P^{(\Gamma - 1)/\Gamma} - P_0^{(\Gamma - 1)/\Gamma}) = -K^{-1/\Gamma}g z
\end{equation}
Since we introduced the concept of our variables at the surface $P_0, \rho_0, T_0$. We can find the expression for $K$ using the polytropic equation of state with our variables at the surface:
\begin{equation}
    P_0 = K\rho_0^{\Gamma}
\end{equation}
Inserting this to our earlier expression we get:
\begin{equation}
    \ab(\frac{\Gamma}{\Gamma - 1})\ab(P^{(\Gamma - 1)/\Gamma} - P_0^{(\Gamma - 1)/\Gamma}) = -P_0^{-1/\Gamma}\rho_0g z
\end{equation}
\begin{equation}
    P^{(\Gamma - 1)/\Gamma} - P_0^{(\Gamma - 1)/\Gamma} = -\ab(\frac{\Gamma-1}{\Gamma})P_0^{-1/\Gamma}\rho_0g z
\end{equation}
\begin{equation}
    P^{(\Gamma - 1)/\Gamma} =P_0^{(\Gamma - 1)/\Gamma} -\ab(\frac{\Gamma-1}{\Gamma})P_0^{-1/\Gamma}\rho_0g z
\end{equation}
\begin{equation}
    P^{(\Gamma - 1)/\Gamma} = P_0^{(\Gamma - 1)/\Gamma}\ab(1-\ab(\frac{\Gamma-1}{\Gamma})P_0^{-1}\rho_0g z)
\end{equation}
\begin{equation}
    P = P_0\ab(1-\ab(\frac{\Gamma-1}{\Gamma})\frac{\rho_0}{P_0}g z)^{\Gamma/(\Gamma-1)}
\end{equation}
\begin{equation}
    \boxed{P = P_0\ab(1-\ab(\frac{\Gamma-1}{\Gamma})\frac{\mu m_H}{kT_0}g z)^{\Gamma/(\Gamma-1)}}
\end{equation}
This serves as the base for our expression of the pressure. Let's go to the temperature for now. We have the ideal gas law:
\begin{equation}
    P = \frac{\rho k T}{\mu m_H}
\end{equation}
\begin{equation}
    \rho = \frac{P}{T}\frac{\mu m_H}{k}
\end{equation}
Let's take a ratio of pressure and the ground pressure in the polytropic equation of state:
\begin{equation}
    P/P_0 = \frac{K}{K}\ab(\frac{\rho}{\rho_0})^{\Gamma}
\end{equation}
Inserting our expression for the density from above:
\begin{equation}
    P/P_0 = \ab(\frac{P}{T}\frac{T_0}{P_0})^\Gamma
\end{equation}
Putting all constants on one side we get:
\begin{equation}
    P^{1-\Gamma}T^{\Gamma} = const
\end{equation}
Where we recover the polytropic equation of state for temperature. We then get:
\begin{equation}
    \frac{T}{T_0} = \ab(\frac{P}{P_0})^{(\Gamma - 1)/\Gamma}
\end{equation}
Since we have our expression for $P$, we get:
\begin{equation}
    T = T_0\ab(\ab(1-\ab(\frac{\Gamma-1}{\Gamma})\frac{\mu m_H}{kT_0}g z)^{\Gamma/(\Gamma-1)})^{(\Gamma - 1)/\Gamma}
\end{equation}
\begin{equation}
    \boxed{T = T_0 - \ab(\frac{\Gamma-1}{\Gamma})\frac{\mu m_H}{k}g z}
\end{equation}
Simply use the polytropic relation again to recover the density profile:
\begin{equation}
    P = K\rho^\Gamma
\end{equation}
\begin{equation}
    \boxed{\rho/\rho_0 = \ab(\frac{P}{P_0})^{1/\Gamma}}
\end{equation}
To find the temperature at mount everest. We can plug in the value of $z$ in the function for $T$, and use a $\Gamma = 7/5$ and $\mu=28$ since we consider most of the atmosphere to be diatomic nitrogen.
\begin{equation}
    T(z = \qty{8848}{m}) = \qty{300}{K} - \ab(\frac{(7/5)-1}{7/5})\frac{28 m_H}{k}(\qty{9.81}{m/s^2})(\qty{8848}{m})
\end{equation}
Using $m_H = \qty{1.67e-27}{kg}, k = \qty{1.38e-23}{[SI]}$ we get:
\begin{equation}
    \boxed{T(z = \qty{8848}{m}) = \qty{216}{K}}
\end{equation}
We can use the polytropic equation of state for temperature to retrieve the ratio of the pressure at z and the ground:
\begin{equation}
    \frac{P}{P_0} = \ab(\frac{T}{T_0})^{\Gamma/(\Gamma-1)}
\end{equation}
\begin{equation}
    \frac{P}{P_0} = \ab(\frac{\qty{216}{K}}{\qty{300}{K}})^{(7/5)/(7/5-1)} = 0.3 
\end{equation}
\begin{equation}
    \boxed{P(z=\qty{8848}{m}) = 0.3P_0}
\end{equation}
The pressure at the top of mount everest is about $30\%$ of the pressure at the ground.